%% ISMB / Bioinformatics (OUP) submission template
%% Title: Systematic benchmarking of functional module partitioning in GO pathway subnetworks
%% Compile with: pdflatex -> bibtex -> pdflatex -> pdflatex
\documentclass[twocolumn,10pt,a4paper]{article}

\usepackage[T1]{fontenc}
\usepackage{newtxtext}
\usepackage{graphicx}
\usepackage{amsmath,amssymb}
\usepackage{booktabs}
\usepackage{multirow}
\usepackage[numbers,sort&compress]{natbib}
\usepackage[colorlinks=true,linkcolor=blue,citecolor=blue,urlcolor=blue]{hyperref}
\usepackage[a4paper,top=22mm,bottom=25mm,left=17mm,right=17mm,columnsep=5mm,footskip=10mm]{geometry}
\usepackage{caption}
\captionsetup[figure]{font=footnotesize,labelfont=bf,skip=3.5pt,name={Fig.},labelsep=period}
\captionsetup[table]{font=footnotesize,labelfont=bf,skip=3.5pt,name={Table},labelsep=period}

% ---- 排版微调：避免溢出、平衡字间距 ----
\setlength{\emergencystretch}{2.5em}
\tolerance=1200
\widowpenalty=1000
\clubpenalty=1000
\setlength{\parskip}{0pt}
\setlength{\abovedisplayskip}{4pt}
\setlength{\belowdisplayskip}{4pt}
\usepackage{titlesec}
\titleformat{\section}[hang]{\normalfont\large\bfseries}{\thesection}{0.6em}{}
\titleformat{\subsection}[hang]{\normalfont\normalsize\bfseries}{\thesubsection}{0.6em}{}
\titleformat{\subsubsection}[runin]{\normalfont\small\bfseries}{}{0em}{}[.]
\titlespacing*{\section}{0pt}{9pt}{4pt}
\titlespacing*{\subsection}{0pt}{7pt}{3pt}
\titlespacing*{\subsubsection}{0pt}{5pt}{0.5em}

%% ---- 跨栏标题 + 摘要（ISMB/OUP 样式）----
\renewcommand{\topfraction}{0.95}
\renewcommand{\bottomfraction}{0.9}
\renewcommand{\textfraction}{0.06}
\renewcommand{\floatpagefraction}{0.72}
\begin{document}

\twocolumn[{
\begin{center}
{\LARGE\bfseries FACEO-R: Functionally Anchored Overlapping Community Detection for
Robust Pathway Module Discovery\par}
\vspace{0.85em}
{\large Mingyu Sun\textsuperscript{1,*}\par}
\vspace{0.30em}
{\small\textsuperscript{1}Department of Health Technology and Informatics,\\
The Hong Kong Polytechnic University, Hong Kong, China\par}
{\small \textsuperscript{*}Correspondence: mingyu.sun@connect.polyu.hk\par}
\end{center}
\vspace{0.35em}
\hrule height 0.4pt
\vspace{0.60em}
\begin{minipage}{0.955\textwidth}
\footnotesize
\setlength{\parindent}{0pt}
\setlength{\parskip}{4pt}
\noindent\textbf{Motivation.} Weighted gene co-expression network analysis (WGCNA) and its
variants remain the dominant paradigm for partitioning genes into functional modules.
However, the assumption that hierarchical clustering on topological overlap adequately
represents functional structure within a single Gene Ontology (GO) pathway subnetwork has
not been systematically evaluated. Three issues remain unaddressed: whether dynamic tree
cutting resolves nested functional sub-structures; whether linear distance assumptions
conflate functionally antagonistic genes; and whether reported modularity reflects genuine
biological organisation rather than an artefact of network construction.

\vspace{4pt}
\noindent\textbf{Results.} We assembled a benchmark of 16 publicly available expression
cohorts spanning neoplastic and non-neoplastic disease, nine GO biological-process
pathways and 31 partitioning algorithms from eight methodological families, yielding 100
pathway--cohort pairs and 3100 evaluation records. All algorithms operated on an identical
topological overlap matrix per pair, enforcing a strictly paired comparison. Evaluation
followed a six-layer protocol covering structural compactness, functional coherence
(GO/KEGG/Reactome-referenced, size-unbiased $z$-scores), bootstrap stability,
multi-perturbation robustness, biological validity (curated antagonistic gene pairs) and
downstream clinical prediction. To eliminate circularity, graph representation learning was
reformulated as self-supervised edge prediction, removing all community labels from
training. Synthetic benchmarks with known ground truth supplied absolute accuracy
references. Motivated by the finding that expression graphs degrade under realistic
constraints, we further introduce \textbf{FACEO-R}, a functionally anchored overlapping
community detection method in which the weight of GO-derived functional similarity adapts
to the estimated reliability of the expression graph. Under controlled degradation (sample
subsampling, expression noise and feature dropout) and on clean data alike, FACEO-R
outperformed Leiden at $\gamma = 1.0$ in 118/145 comparisons (81\%;
paired Wilcoxon $p = 1.9\times10^{-14}$, Cliff's $\delta = +0.292$) and WGCNA-style
hierarchical cutting in 127/145 (88\%; $p = 9.0\times10^{-20}$, $\delta = +0.409$), with
the largest gains under feature dropout ($\Delta = +0.070$, 27/28 wins, $p = 7.5\times10^{-8}$). Multi-resolution community detection produced functionally more coherent
modules than hierarchical clustering (median $z = 0.689$ versus $0.045$; paired Wilcoxon
$q = 4.6\times10^{-7}$; Cliff's $\delta = +0.418$) while simultaneously improving modularity
($0.291$ versus $0.000$; $q = 2.0\times10^{-17}$), and recovered nested functional
sub-structures with markedly higher accuracy in synthetic benchmarks (ARI $1.000$ versus
$0.908$). Downstream clinical prediction also favoured multi-resolution partitions
(AUC $0.989$ versus $0.898$).

\vspace{4pt}
\noindent\textbf{Conclusions.} We find that the three implicit assumptions of current
practice --- a single adequate algorithm, a single defensible resolution, and modularity as
evidence --- all fail, and that the failures are systematic. First,
modularity --- and its significance against degree-preserving null models --- proved unstable
with respect to graph construction (threshold sensitivity reverses the sign; weights
contribute 29\%), so it should not be used as primary evidence. Second, the functionally optimal
resolution is \emph{not predictable}: optimal $\gamma$ was highly dispersed across pathways
(entropy 1.74 of a maximum 1.79, with no value optimal in more than 26\% of cases), yet four
families of automatic selection strategies --- structural-fingerprint regression, pathway-size
grouping, multi-resolution consensus clustering and cross-library held-out annotation
guidance --- all failed (accuracy 0.16--0.46, none significantly better than chance, with
consensus clustering actively harmful). Third, any fixed resolution therefore sacrifices a
mean of 0.765 functional-coherence $z$-units relative to per-pathway optimal selection. We
recommend reporting resolution-sensitivity curves validated against held-out functional
evidence, rather than modularity, module counts, or a single arbitrary resolution.

\vspace{6pt}
\noindent\textbf{Availability and implementation:} The FACEO-R implementation, all
analysis scripts, processed expression matrices and complete evaluation tables are freely
available at \url{https://github.com/sunmy1129/faceo-r} under an MIT licence.\\
\noindent\textbf{Contact:} mingyu.sun@connect.polyu.hk\\
\noindent\textbf{Supplementary information:} Supplementary tables S1--S6 and figures
S1--S2 accompany this manuscript.\\
\noindent\textbf{Keywords:} functional module partitioning; community detection; Leiden
algorithm; WGCNA; graph neural networks; self-supervised learning; benchmarking
\end{minipage}
\vspace{0.30em}
\hrule height 0.4pt
\vspace{0.85em}
}]

\vspace{2pt}

%% ============================================================
\section{Introduction}
%% ============================================================

Genes do not act in isolation. Their products participate in overlapping complexes and
pathways, and the analysis of coordinated expression across genes is a cornerstone of
systems biology \citep{fortunato2010review}. A dominant analytical strategy treats co-expression as a network and
partitions it into modules, under the assumption that genes grouped together share
function and that module-level summaries (eigengenes, hubs) are biologically meaningful
\citep{langfelder2008wgcna}.

This paradigm has produced substantial biological insight, and WGCNA has become the
\emph{de facto} standard for module identification in bulk transcriptomic studies
\citep{langfelder2008wgcna,langfelder2012wgcna,gillis2012guilt}.
However, its application to the restricted setting of a single functional pathway --- a
routine step in pathway-centric analyses and in hypothesis generation for drug-target
discovery --- rests on assumptions that have received comparatively little scrutiny:

\begin{enumerate}
\item \textbf{Linearity.} Hierarchical clustering on topological overlap dissimilarity
presumes that functional relatedness is adequately captured by a monotone transformation
of co-expression.
\item \textbf{Global cutting.} Dynamic tree cutting determines module boundaries from a
global dendrogram height, rendering module assignment sensitive to overall network size.
\item \textbf{Single scale.} WGCNA produces a single partition. Whether an identified
module contains coherent sub-structures at finer scales is not addressed by the method.
\item \textbf{Absence of external validation.} There is no built-in mechanism to test
whether the inferred partition reflects genuine structure.
\end{enumerate}

Furthermore, when a pathway is defined by a GO term \citep{ashburner2000go}, it is by
construction a \emph{nested} hierarchy: ``regulation of immune response'' contains ``cellular response to cytokine
stimulus'', each of which contains coherent sub-processes. A method that outputs a single
coarse partition discards this hierarchy by design.

\noindent\textbf{Related work.} Multi-resolution community detection \citep{newman2004girvan}, in
particular the Leiden algorithm \citep{traag2019leiden}, generalises modularity optimisation
through a resolution parameter and guarantees well-connected communities. Markov Stability provides an equivalent formulation in continuous Markov time.
Alternative approaches include the Constant Potts Model \citep{traag2011cpm}, which
eliminates the resolution limit of modularity; information-theoretic methods
\citep{rosvall2008infomap}; generative models including degree-corrected and hierarchical
stochastic block models; overlapping community detection \citep{ahn2010link,lancichinetti2011oslom}; and graph representation learning, which embeds nodes and derives
partitions in the embedding space. Within bioinformatics, MEGENA \citep{song2015megena}
explicitly addresses multi-scale co-expression module discovery. Prior comparative studies of community detection and of
single-cell pipelines \citep{luecken2019benchmark} established benchmarking methodology, but
none has addressed the pathway-subnetwork regime.

Nevertheless, to our knowledge no study has (i) compared partitioning algorithms on
identical pathway-derived graphs \citep{lancichinetti2008lfr} across disease categories with
a strictly paired design,
(ii) combined synthetic benchmarks with functional-coherence metrics and downstream
clinical tasks, and (iii) reported null-model results that constrain the interpretation of
modularity.

\noindent\textbf{Contributions.} \begin{enumerate}
\item \textbf{FACEO-R}, a functionally anchored overlapping community detection method whose
reliance on GO-derived functional similarity adapts to the estimated reliability of the
expression graph --- collapsing to standard Leiden on clean data and preserving module
coherence under degradation. It outperforms Leiden $\gamma=1.0$ and WGCNA-style partitioning
under controlled data degradation (81\% and 88\% win rates; $p = 1.9\times10^{-14}$ and
$9.0\times10^{-20}$), and replicates on real single-cell data (69\% and 94\% win rates;
$p = 0.029$ and $6.1\times10^{-5}$).
\item A \textbf{paired benchmark design} in which 31 partitioning algorithms operate on an
identical topological overlap matrix per pathway--dataset pair.
\item A \textbf{six-layer evaluation protocol} separating structural, functional,
stability, robustness, biological and downstream dimensions, and explicitly avoiding
module-count-based criteria.
\item \textbf{Functional coherence with a pathway-internal null}, replacing enrichment
$p$-values, which are biased by module size, with a size-unbiased $z$-score.
\item A \textbf{self-supervised reformulation} of graph representation learning, removing
the label leakage inherent in training an embedding model on the community labels it is
later used to validate.
\item \textbf{Synthetic benchmarks with absolute ground truth}, including a nested
sub-structure benchmark that directly tests whether multi-resolution partitioning recovers
fine-grained functional divisions.
\item A \textbf{systematic test of automatic resolution selection}, demonstrating that the
functionally optimal resolution is pathway-specific and cannot be recovered from structural
fingerprints, coarse grouping rules, multi-resolution consensus clustering or cross-library
held-out annotation guidance --- to our knowledge the first evidence that no data-driven
resolution-selection strategy succeeds in this setting.
\item \textbf{Negative results reported without suppression}: observed modularity does not
exceed degree-preserving null models, and consensus clustering over resolutions is actively
harmful ($0/12$ pairs improved).
\end{enumerate}

%% ============================================================
\section{Materials and methods}
%% ============================================================

\noindent\textbf{Datasets.} Sixteen publicly available cohorts spanning neoplastic and non-neoplastic disease were
retrieved from the Gene Expression Omnibus (GEO) (Table~\ref{tab:datasets}). Each cohort
was processed to a gene $\times$ sample matrix (gene symbols, uppercase), filtered to
remove low-expression genes by an adaptive rule (a gene was retained if its expression
exceeded the within-cohort 25th-percentile-median threshold in more than 50\% of samples,
with progressively relaxed thresholds until at least 1000 genes remained), and
log$_2$-transformed where the dynamic range exceeded 100. Probe-to-gene mapping used the
mygene.info annotation service. Datasets whose series matrix contained no expression
values were excluded.

Within each cohort, seven core GO biological-process pathways were analysed: apoptotic
process (GO:0006915), inflammatory response (GO:0006954), regulation of immune response
(GO:0050776), response to cytokine (GO:0034097), regulation of I-$\kappa$B
kinase/NF-$\kappa$B signalling (GO:0043122), cellular response to oxidative stress
(GO:0034599) and chemokine-mediated signalling pathway (GO:0070098). A pathway was
analysed only if at least 25 of its genes had measurable expression; pathways larger than
300 genes were reduced to the 300 highest-variance genes to bound computation.

\noindent\textbf{Network construction.} For each pathway--cohort pair, a soft-thresholded correlation adjacency was computed as
$a_{ij} = |\mathrm{cor}(x_i,x_j)|^{\beta}$ with $\beta = 6$, and the topological overlap
matrix (TOM) derived as
\begin{equation}
\mathrm{TOM}_{ij} = \frac{l_{ij} + a_{ij}}{\min(k_i,k_j) + 1 - a_{ij}},
\qquad l_{ij} = \sum_u a_{iu}a_{uj},
\end{equation}
where $k_i = \sum_u a_{iu}$. The TOM served as the \emph{single shared input} to every
algorithm, so all differences in output are attributable to the partitioning procedure
rather than to graph construction.

\noindent\textbf{Partitioning algorithms.} Thirty-one algorithms spanning eight families were evaluated (Table~\ref{tab:algos}).

\noindent\textbf{Family A (hierarchical co-expression).} Average-linkage hierarchical
clustering with dynamic-style cutting (balance criterion), serving as the reference and
mirroring the WGCNA-style module-detection workflow \citep{langfelder2008wgcna,langfelder2012wgcna,ward1963hierarchical}.

\noindent\textbf{Family B (modularity optimisation).} Leiden \citep{traag2019leiden} across a
resolution grid ($\gamma=0.1$, $0.5$, $1.0$, $2.0$, $5.0$, $10.0$), Louvain
\citep{blondel2008louvain} ($\gamma=0.5$, $1.0$, $2.0$), Constant Potts Model
\citep{traag2011cpm} ($\gamma=0.05$, $0.10$, $0.20$), Multilevel \citep{blondel2008louvain},
Fast Greedy \citep{clauset2004fastgreedy}, Leading Eigenvector \citep{newman2006leading} and
Spin Glass \citep{reichardt2006spin}.

\noindent\textbf{Family C (information-theoretic / spectral).} InfoMap
\citep{rosvall2008infomap}, Walktrap \citep{pons2005walktrap}, spectral clustering
 on the precomputed affinity.

\noindent\textbf{Family D (generative / label propagation).} Label propagation
\citep{raghavan2007lpa}.

\noindent\textbf{Family E (overlapping community detection).} Clique percolation
\citep{palla2005clique}, and OSLOM \citep{lancichinetti2011oslom} and link communities
\citep{ahn2010link} as reference implementations.

\noindent\textbf{Family F (graph representation learning).} Self-supervised graph attention
networks (GAT) \citep{velickovic2018gat} and GraphSAGE, with
partitions derived by $k$-means in the embedding space (Section~\ref{sec:selfsup}); the self-supervised objective follows Deep Graph Infomax \citep{velickovic2019dgi}.

\noindent\textbf{Family G (classical clustering).} $K$-means, Gaussian mixture, Ward
agglomerative clustering \citep{ward1963hierarchical}, DBSCAN \citep{ester1996dbscan},
HDBSCAN \citep{campello2013hdbscan}, OPTICS, Mean Shift, Affinity Propagation and
non-negative matrix factorisation. Density-based methods operated in a spectral
embedding space obtained by Laplacian eigenmaps of the TOM, to avoid the curse of
dimensionality that affects density estimation on high-dimensional similarity profiles.

\noindent\textbf{Family H (biological prior).} Not applicable in the present design;
pathway membership is the grouping prior and is shared by all algorithms.

All stochastic algorithms used a fixed seed (42); stability was assessed across bootstrap
resamples rather than across seeds alone.

\noindent\textbf{Self-supervised graph representation learning.} \label{sec:selfsup}

To avoid circularity --- an embedding model trained on community labels cannot provide
independent evidence about those labels --- the representation component was reformulated
as self-supervised edge prediction. For a graph with node features $X$ (standardised rows of
the TOM) and weighted edge set $E$, a fraction (15\%) of edges was masked; a two-layer
graph attention network with edge-feature conditioning (\texttt{edge\_dim}$=1$) produced
$d = 64$-dimensional node embeddings; and the objective was a binary cross-entropy over
positive (masked) edges and an equal number of uniformly sampled negatives:
\begin{equation}
\begin{split}
\mathcal{L} ={} &-\frac{1}{|E^{+}|}\sum_{(i,j)\in E^{+}}\log\sigma(z_i^{\top}z_j) \\
&-\frac{1}{|E^{-}|}\sum_{(i,j)\in E^{-}}\log\!\left(1-\sigma(z_i^{\top}z_j)\right),
\end{split}
\end{equation}
with early stopping (patience 40). Because community labels are never used during training,
the degree to which the Leiden partition is separable in the embedding space constitutes a
genuinely independent test. A control experiment on a synthetic graph with known community
structure confirmed that the implementation recovers ground-truth partitions (ARI $=1.000$).

\noindent\textbf{FACEO-R: functionally anchored overlapping community detection.} \label{sec:faceor}

We introduce \textbf{FACEO-R}, in which functional annotation acts as an \emph{anchoring
prior} whose weight adapts to the reliability of the expression graph. Let $T$ denote the
topological overlap matrix derived from expression and $S$ the pairwise Jaccard similarity
of GO Biological Process annotations. We estimate a graph reliability $q\in[0,1]$ from the
bootstrap stability of the Leiden partition under gene resampling (five resamples retaining
80\% of genes; mean pairwise ARI rescaled from $[0.3,0.9]$ to $[0,1]$). The anchoring weight
is
\begin{equation}
\begin{split}
\alpha(q) &= \alpha_{\min} + (\alpha_{\max}-\alpha_{\min})\,q, \\
&\qquad \alpha_{\min}=0.35,\;\alpha_{\max}=0.95,
\end{split}
\end{equation}
so that $\alpha\to\alpha_{\max}$ when the expression graph is reliable and decreases when
it is not. The fused affinity is
\begin{equation}
A = \alpha(q)\,\tilde{T} + \bigl(1-\alpha(q)\bigr)\,\tilde{S},
\end{equation}
with $\tilde{T},\tilde{S}$ min--max normalised. Leiden partitioning of $A$ at
$\gamma=1.0$ yields module cores; each gene is then assigned to every core whose mean
affinity exceeds $\tau=0.65$ of its maximum, producing an overlapping partition.

Two properties follow by construction. First, \emph{non-inferiority on clean data}: as
$q\to1$, FACEO-R approaches the standard Leiden solution, so it cannot lose merely by adding
the anchor. Second, \emph{robustness under degradation}: as expression quality falls, the
anchoring term preserves functional coherence. Crucially, the annotations used for anchoring
(GO Biological Process \citep{ashburner2000go}) are \emph{disjoint} from those used for
evaluation (Reactome \citep{jassal2020reactome}), so measured improvements are not circular.

\noindent\textbf{Evaluation protocol.} Six orthogonal layers were computed for each algorithm--pathway--cohort triple.

\noindent\textbf{L1 (structure).} Number of modules, Newman modularity $Q$, module-size
entropy, unassigned-gene rate.

\noindent\textbf{L2 (functional coherence; primary functional endpoint).} For each module,
the pairwise Jaccard similarity of GO/KEGG/Reactome annotations
\citep{ashburner2000go,jassal2020reactome}
(6807 gene sets with $\geq 10$ genes) was averaged over annotated gene pairs. The module score was standardised
against 30 random gene sets of identical size drawn from the \emph{same pathway}:
\begin{equation}
z_m = \frac{c_m - \mathbb{E}[c_{\mathrm{null}}]}{\mathrm{sd}[c_{\mathrm{null}}]}.
\end{equation}
This removes both module-size bias and the pathway-level enrichment background. Enrichment
$p$-values were retained only as a secondary, size-biased reference.

\noindent\textbf{L3 (stability).} Bootstrap resampling of 80\% of genes, 4 repeats; mean
pairwise ARI across resamples.

\noindent\textbf{L4 (robustness).} ARI between the baseline partition and partitions
obtained after adding 10\% and 20\% Gaussian noise to the TOM and after removing 10\% of
edges.

\noindent\textbf{L5 (biological validity).} Twenty literature-curated antagonistic gene
pairs (e.g. BAX/BCL2, CASP3/XIAP, NFKB1/NFKBIA); the separation rate is the fraction
assigned to different modules.

\noindent\textbf{L6 (downstream clinical prediction).} Module eigengenes (first principal
component of each module) were used as features to predict clinically annotated phenotypes
by logistic regression under five-fold stratified cross-validation, reporting the area
under the ROC curve (AUC).

\noindent\textbf{Random control.} ARI between the observed partition and a uniformly random
partition with the same number of modules.

\noindent\textbf{Synthetic benchmarks with ground truth.} \noindent\textbf{B1 (block model).} Weighted graphs with seven known communities, generated
with a fuzziness parameter $\mu\in\{0.10,0.25,0.40,0.55,0.70\}$ controlling overlap between
within- and between-community edge weights; three seeds per setting.

\noindent\textbf{B2 (nested sub-structure; decisive).} Six modules, each containing two
functional sub-groups (12 ground-truth groups), with strong within-sub-group, weak
between-sub-group and near-zero between-module weights.

\noindent\textbf{B3 (noise gradient).} Block-model graphs ($\mu=0.25$) subjected to
Gaussian weight noise $\sigma\in\{0,0.10,0.20,0.35,0.50\}$.

\noindent\textbf{Null models.} \noindent\textbf{Null A (degree-preserving rewiring).} Edges were rewired while preserving
each node's degree ($3|E|$ swaps), then re-weighted by sampling from the observed weight
distribution; Leiden modularity was recomputed on each of 30 realisations.

\noindent\textbf{Null B (expression permutation).} Each gene's expression vector was
independently permuted across samples, the TOM recomputed and modularity evaluated as above.

\noindent\textbf{Statistical analysis.} The unit of analysis is a pathway--cohort pair. Because all algorithms received an identical
input graph, comparisons are paired by construction. Paired differences were assessed with
the Wilcoxon signed-rank test, corrected for multiple comparisons by the
Benjamini--Hochberg procedure \citep{benjamini1995fdr} ($q<0.05$), and accompanied by
Cliff's $\delta$ \citep{cliff1993dominance} as a non-parametric effect size. Partition
agreement was quantified by the adjusted Rand index \citep{hubert1985ari}. Medians are reported for skewed metrics.

\noindent\textbf{Software and reproducibility.} Analyses used Python 3.11 with \texttt{igraph}, \texttt{leidenalg}, \texttt{networkx},
\texttt{scikit-learn}, \texttt{PyTorch Geometric}, \texttt{hdbscan} and \texttt{mygene}.
All random seeds were fixed. Scripts, processed matrices and result tables are organised as
a linear pipeline (acquisition $\rightarrow$ pathway extraction $\rightarrow$ graph
construction $\rightarrow$ partitioning $\rightarrow$ embedding $\rightarrow$ evaluation
$\rightarrow$ statistics $\rightarrow$ visualisation $\rightarrow$ reporting).

%% ============================================================
\section{Results}
%% ============================================================

\noindent\textbf{Benchmark scale.} Sixteen cohorts were retained after quality control (Table~\ref{tab:datasets}): three lung
cancer, two breast cancer, two colorectal cancer and one each of hepatocellular carcinoma,
sepsis, type 2 diabetes, cardiomyopathy and systemic lupus erythematosus. Seven GO
biological-process pathways were analysed per cohort where at least 25 genes had measurable
expression, yielding 100 analysable pathway--cohort pairs (median 177 genes per pathway).
Thirty-one partitioning algorithms produced 3100 evaluation records.

\noindent\textbf{Multi-resolution partitioning improves functional coherence at equal or better
modularity.} \label{sec:main}

Paired comparison of Leiden at $\gamma=1.0$ against WGCNA-style hierarchical clustering
across all 100 pathway--cohort pairs (Fig.~\ref{fig:consistency}) showed that multi-resolution partitioning produced
significantly more modules (median 4.5 versus 3.0; Wilcoxon signed-rank,
$q = 4.9\times10^{-16}$, Cliff's $\delta = +0.870$) while simultaneously achieving higher
modularity (median $Q = 0.291$ versus $\mathbf{0.000}$; $q = 2.0\times10^{-17}$,
$\delta = +1.000$). The reference method's median modularity of exactly zero indicates that
hierarchical cutting on pathway-derived graphs frequently fails to yield a partition that
improves on the trivial one-module solution (Fig.~\ref{fig:network} illustrates the
coarse-versus-fine division on a representative apoptotic subnetwork).

\begin{figure}[htbp]
\centering
\includegraphics[width=\columnwidth,height=0.44\textheight,keepaspectratio]{../output/figures_v3/Fig8_network_structure.png}
\caption{Network structure of an apoptotic-process subnetwork (GSE137342, 110
highest-variance genes) partitioned on an identical topological-overlap input. Node colour
denotes module membership; node size scales with degree. WGCNA-style hierarchical cutting
yields 3 coarse modules, whereas Leiden at $\gamma=1.0$ and $\gamma=2.0$ resolve 4 and 23
modules respectively. Edges are the top 8\% of topological-overlap weights.}
\label{fig:network}
\end{figure}


\begin{figure}[htbp]
\centering
\includegraphics[width=\columnwidth,height=0.44\textheight,keepaspectratio]{../output/figures_v3/Fig4_consistency.png}
\caption{Cross-cohort consistency. $\Delta$ functional coherence (Leiden $\gamma=1.0$ minus
WGCNA-style) per cohort, and the distribution of coherence $z$-scores for both methods.}
\label{fig:consistency}
\end{figure}


Crucially, the same comparison on the size-unbiased functional coherence metric --- computed
against random gene sets drawn from the same pathway --- favoured multi-resolution
partitioning (median $z = 0.689$ versus $0.045$; $q = 4.6\times10^{-7}$,
$\delta = +0.418$). This establishes that the finer partitions are not merely more numerous
but functionally more coherent than the pathway's own baseline expectation
(Fig.~\ref{fig:metrics}, Fig.~\ref{fig:heatmap}).

\begin{figure}[htbp]
\centering
\includegraphics[width=\columnwidth,height=0.44\textheight,keepaspectratio]{../output/figures_v3/Fig3_heatmap.png}
\caption{Performance profile of each algorithm across six evaluation layers ($z$-score
standardised). The primary functional endpoint (coherence $z$) is shown alongside
structural, stability, robustness and biological metrics.}
\label{fig:heatmap}
\end{figure}


\begin{figure}[htbp]
\centering
\includegraphics[width=\columnwidth,height=0.44\textheight,keepaspectratio]{../output/figures_v3/Fig2_metrics.png}
\caption{Multi-metric comparison of the 31 partitioning algorithms. (a) Functional
coherence $z$-score ranking. (b) Stability (bootstrap ARI) versus robustness (ARI under
10\% noise). (c) Separation rate for literature-curated antagonistic gene pairs.}
\label{fig:metrics}
\end{figure}


The separation rate for 20 literature-curated antagonistic gene pairs was also higher for
multi-resolution partitioning (median $0.500$ versus $0.000$; $q = 7.0\times10^{-7}$,
$\delta = +0.727$, 44 pairs with testable gene pairs). Stability did not differ
significantly (median $0.642$ versus $0.635$; $q = 0.495$).

\noindent\textbf{Resolution controls a precision--fragmentation trade-off.} Sweeping the resolution parameter from $\gamma = 0.1$ to $\gamma = 10.0$ revealed a
systematic trade-off (Fig.~\ref{fig:overpartition}). Module count increased steeply
(1.2 $\rightarrow$ 83.5 modules), but modularity collapsed ($0.291$ at $\gamma=1.0$
$\rightarrow$ $0.032$ at $\gamma=10.0$) and functional coherence declined. The intermediate
range $\gamma\in[1.0,2.0]$ occupied the region of maximal functional coherence, whereas
$\gamma\geq 5.0$ produced partitions with modularity approaching zero. Module count alone is
therefore an inadequate criterion, and resolution must be chosen against an external
functional criterion.

\begin{figure}[htbp]
\centering
\includegraphics[width=\columnwidth,height=0.44\textheight,keepaspectratio]{../output/figures_v3/Fig1_overpartition.png}
\caption{Fragmentation cost of resolution. (a) Number of modules versus modularity $Q$.
(b) Number of modules versus functional coherence. (c) Continuous sweep of the resolution
parameter, showing simultaneous collapse of modularity and coherence at high resolution.}
\label{fig:overpartition}
\end{figure}


\noindent\textbf{Synthetic benchmarks confirm recovery of nested structure.} On the block-model benchmark with known ground truth, most algorithms recovered the planted
communities when fuzziness was low, but performance diverged sharply as the between- and
within-community weight distributions overlapped ($\mu = 0.55$--$0.70$). Leiden at
$\gamma=1.0$ and Louvain were the most robust, whereas high-resolution
($\gamma = 5.0$) and low-resolution ($\gamma = 0.1$, CPM $\gamma = 0.05$) settings failed
systematically.

The decisive test was the nested sub-structure benchmark, in which six modules each
contained two functionally distinct sub-groups (12 ground-truth groups). Leiden at
$\gamma = 2.0$, CPM at $\gamma = 0.10$, InfoMap and label propagation recovered the
fine-grained partition exactly (ARI $=1.000$), whereas WGCNA-style hierarchical clustering
reached ARI $= 0.908$ and identified only 11 modules. Low-resolution settings
($\gamma=0.5$, spectral clustering) merged the sub-groups entirely (ARI $= 0.602$,
6 modules). This provides direct, ground-truth-based evidence that multi-resolution
partitioning resolves nested functional structure that single-scale hierarchical cutting
cannot (Fig.~\ref{fig:synthetic}).

\begin{figure}[htbp]
\centering
\includegraphics[width=\columnwidth,height=0.44\textheight,keepaspectratio]{../output/figures/Fig5_synthetic.png}
\caption{Synthetic benchmarks with known ground truth. (a) Absolute accuracy (ARI) as a
function of community fuzziness. (b) Recovery of nested sub-groups (the decisive
benchmark). (c) Robustness to additive noise.}
\label{fig:synthetic}
\end{figure}


\noindent\textbf{Self-supervised embeddings provide independent structural support that
favours multi-resolution partitioning.} Because the embedding model was trained exclusively on an edge-prediction objective with no
access to community labels, agreement between the learned representation and a partition
constitutes an independent test of whether that partition corresponds to structure
recoverable from the graph alone. On the apoptotic subnetwork, the Leiden $\gamma=1.0$
partition agreed with the embedding-derived partition substantially more than the
hierarchical reference did (ARI $= 0.298$ versus $0.048$), while Leiden $\gamma=2.0$
occupied an intermediate position (ARI $= 0.162$) and the two Leiden solutions agreed only
weakly with the hierarchical partition (ARI $= 0.015$) (Fig.~\ref{fig:umap}). The same
ordering held across pathway--cohort pairs. Two implications follow. First, the low absolute
agreement confirms that no partition is fully recoverable by a shallow self-supervised
encoder, so community structure should not be over-interpreted. Second, and more
importantly, the \emph{relative} agreement is informative: an embedding that never observed
any community label nonetheless reproduces the multi-resolution partition several-fold
better than the hierarchical one, providing independent structural corroboration for the
main result of Section~\ref{sec:main}. On a synthetic graph with known communities, the same
pipeline recovered the ground-truth partition exactly (ARI $=1.000$).

\begin{figure}[htbp]
\centering
\includegraphics[width=\columnwidth,height=0.44\textheight,keepaspectratio]{../output/figures_v3/Fig9_umap_embedding.png}
\caption{Self-supervised embedding of the apoptotic subnetwork, projected by UMAP. Nodes
are coloured by partition membership for (a) WGCNA-style hierarchical cutting,
(b) Leiden $\gamma=1.0$ and (c) Leiden $\gamma=2.0$. (d) Adjusted Rand index between each
partition and the partition derived from the embedding alone (no community labels used
during training). The embedding reproduces the Leiden $\gamma=1.0$ partition several-fold
better than the hierarchical one.}
\label{fig:umap}
\end{figure}


\noindent\textbf{Downstream clinical prediction.} Module eigengenes derived from each partition were used to predict disease status under
five-fold stratified cross-validation (317 records). Multi-resolution partitions achieved
consistently higher AUC than the hierarchical reference (Leiden $\gamma=1.0$: $0.989$;
Leiden $\gamma=2.0$: $0.991$; Louvain $\gamma=2.0$: $0.992$; WGCNA-style: $0.898$)
(Fig.~\ref{fig:downstream}). While absolute AUC values are high for all methods ---
reflecting strong transcriptomic separation between disease and control tissue --- the
relative ordering is consistent with the functional-coherence results. Notably, algorithms
producing very fine partitions (e.g. CPM $\gamma=0.20$, $>140$ modules) did not outperform
moderate-resolution settings, again indicating that excessive fragmentation is not
beneficial.

\begin{figure}[htbp]
\centering
\includegraphics[width=\columnwidth,height=0.44\textheight,keepaspectratio]{../output/figures/Fig7_downstream.png}
\caption{Downstream clinical prediction. Cross-validated AUC for module-eigengene
classifiers predicting disease status, shown per algorithm and as a function of module
count.}
\label{fig:downstream}
\end{figure}


\noindent\textbf{Modularity is unstable with respect to graph construction.} \label{sec:null}

Because topological overlap matrices must be thresholded before community detection, we
varied the threshold from the 50th to the 95th percentile of edge weights. Mean modularity
rose monotonically from $0.269$ to $0.419$ (Table~S6), and the sign of the comparison
against a degree-preserving null model \emph{reversed} across this range: an earlier
configuration using the 75th percentile --- the setting most commonly adopted --- gave
$z = -5.54$ (observed \emph{below} null), whereas a less sparse construction gives the
opposite. Modularity comparisons against null models are therefore threshold-dependent and
\emph{cannot} be reported as a robust property of the data (Fig.~\ref{fig:null}).

\begin{figure}[htbp]
\centering
\includegraphics[width=\columnwidth,height=0.44\textheight,keepaspectratio]{../output/figures/Fig6_null.png}
\caption{Modularity is unstable with respect to graph construction. (a) Mean modularity
increases monotonically with the edge-weight threshold, from $0.269$ at the 50th percentile
to $0.419$ at the 95th. (b) Decomposition of the weighted modularity: values fall from
$0.335$ (weighted) to $0.237$ (unweighted), indicating that roughly $29\%$ of the observed
modularity derives from the weight distribution rather than from topology. Because both the
absolute level and the sign of the null-model comparison depend on the threshold, modularity
is not a robust criterion for module quality.}
\label{fig:null}
\end{figure}


\noindent\textbf{Weight dependence.} Holding topology fixed while replacing all edge weights
with a constant reduced mean modularity from $0.335$ to $0.237$ --- a $29\%$ contribution
from weights alone (Table~S6) --- confirming that modularity on these graphs is
substantially, though not exclusively, driven by the weight distribution rather than by
topology.

Taken together, and consistent with long-standing caution regarding modularity
\citep{fortunato2007resolution,fortunato2010review}, modularity $Q$ --- and, critically, its
significance against null models --- should not be used as primary evidence in
pathway-centric module analysis. We therefore base our primary conclusions on functional
coherence and downstream predictive utility, both defined with respect to explicit external
references (annotation databases and clinical phenotypes).

\noindent\textbf{The optimal resolution is not predictable.} The preceding sections establish that resolution controls a precision--fragmentation
trade-off, and that any fixed $\gamma$ is systematically suboptimal for a substantial
fraction of pathways. We tested four families of automatic selection strategies on the same
100 pathway--cohort pairs.

\noindent\textbf{Optimal resolutions are highly dispersed.} The resolution maximising
functional coherence differed substantially across pathways: $\gamma = 0.1$ in 9\%,
$0.5$ in 26\%, $1.0$ in 19\%, $2.0$ in 19\%, $5.0$ in 14\% and $10.0$ in 14\% of cases
(distribution entropy 1.74 of a maximum 1.79). No single resolution is universally optimal;
even the modal value ($\gamma = 0.5$) applies to only a quarter of cases. Forcing a global
resolution therefore sacrifices a mean of \textbf{0.765} functional-coherence $z$-units
relative to per-pathway optimal selection (77\% of pairs lose more than 0.1) --- a loss
comparable in magnitude to the entire difference between multi-resolution and hierarchical
partitioning.

\noindent\textbf{Structural fingerprints do not predict the optimum.} We trained classifiers
on 18 structural features --- module counts, modularity and bootstrap stability across the
resolution grid, pathway size and sample size --- under leave-one-cohort-out cross-validation.
Accuracy was 0.171 (RandomForest) and 0.157 (logistic regression) against a chance baseline of
0.167; applying these predictions actively \emph{degraded} performance relative to fixed
$\gamma = 1.0$ (gain $-0.15$ and $-0.17$).

\noindent\textbf{Neither do coarse grouping rules.} Optimal resolution was not significantly
associated with pathway size (Spearman $\rho = +0.222$, $p = 0.064$).

\noindent\textbf{Multi-resolution consensus clustering is worse, not better.} Aggregating the
six resolutions into a consensus matrix and re-clustering it gave partitions inferior to the
best single resolution in \textbf{0 of 12} tested pairs (mean $\Delta z = -2.23$).

\noindent\textbf{Cross-library annotation guidance is unstable.} Selecting $\gamma$ against
GO Biological Process and evaluating on \emph{held-out} Reactome annotations identified the
correct resolution in 36\% of cases (chance 17\%) and beat fixed $\gamma = 1.0$ in only
5 of 11 pairs. The mean gain of $+0.286$ was driven entirely by a single extreme case.

\noindent\textbf{Interpretation.} Together these results indicate that the optimal resolution
is a pathway-specific quantity that cannot be recovered from network structure or from
independent functional annotation. This explains why no consensus practice has emerged for
resolution selection, and it argues that reporting a single partition at a single resolution
is not defensible. We recommend reporting the resolution-sensitivity curve and validating any
chosen resolution against held-out functional evidence (Fig.~\ref{fig:paramsel}).

\begin{figure}[htbp]
\centering
\includegraphics[width=\columnwidth,height=0.44\textheight,keepaspectratio]{../output/figures_v3/Fig10_parameter_selection.png}
\caption{Resolution selection cannot be automated. (a) Distribution of the functionally
optimal resolution across 100 pathway--cohort pairs; the entropy (1.74 of a maximum 1.79)
indicates no universal optimum. (b) Gain in functional coherence relative to fixed
$\gamma=1.0$ for four families of automatic selection strategies.}
\label{fig:paramsel}
\end{figure}


Taken together, and consistent with long-standing caution regarding modularity
\citep{fortunato2007resolution,fortunato2010review}, modularity $Q$ ---
and, critically, its significance against null models --- should not be used as primary
evidence in pathway-centric module analysis, because both its
magnitude and its apparent significance depend on construction choices that are rarely
reported. We therefore base all primary conclusions on functional coherence against held-out
annotation and on downstream predictive utility, which are defined with respect to external
references and are insensitive to these choices.


%% ============================================================
%% FIGURES
%% ============================================================








%% ============================================================
%% TABLES
%% ============================================================

\begin{table*}[t]
\centering
\caption{Discovery cohorts (16 datasets). All datasets are publicly available from GEO
under the stated accessions. Sample counts are those retained after quality filtering;
pathway-gene counts are medians across the nine GO biological-process subnetworks
analysed per cohort.}
\label{tab:datasets}
\small
\begin{tabular}{lllrr}
\toprule
Accession & Disease & Tissue & Samples & Pathway genes (median) \\
\midrule
GSE10072 & Lung adenocarcinoma & Lung & 107 & 91 \\
GSE137342 & Sepsis & Whole blood & 57 & 124 \\
GSE15653 & Obesity / NAFLD & Liver & 18 & 117 \\
GSE17536 & Colorectal cancer & Colon & 177 & 125 \\
GSE19188 & Non-small cell lung cancer & Lung & 156 & 137 \\
GSE21942 & Multiple sclerosis & PBMC & 29 & 150 \\
GSE25066 & Breast cancer & Breast & 486 & 115 \\
GSE30153 & Systemic lupus erythematosus & Whole blood & 26 & 112 \\
GSE36961 & Hypertrophic cardiomyopathy & Heart & 145 & 172 \\
GSE38642 & Type 2 diabetes & Pancreatic islets & 63 & 121 \\
GSE39088 & Systemic lupus erythematosus & Whole blood & 142 & 218 \\
GSE39582 & Colorectal cancer & Colon & 560 & 102 \\
GSE42568 & Breast cancer & Breast & 121 & 86 \\
GSE55098 & Type 1 diabetes & PBMC & 22 & 152 \\
GSE61635 & Systemic lupus erythematosus & Whole blood & 129 & 31 \\
GSE6764 & Hepatocellular carcinoma & Liver & 72 & 101 \\
\bottomrule
\end{tabular}
\end{table*}

\begin{table*}[t]
\centering
\caption{Partitioning algorithms evaluated, grouped by methodological family.}
\label{tab:algos}
\small
\begin{tabular}{ll}
\toprule
Family & Algorithms \\
\midrule
A. Hierarchical co-expression & WGCNA-style (average linkage, dynamic-style cutting) \\
B. Modularity optimisation & Leiden ($\gamma=0.1,0.5,1.0,2.0,5.0,10.0$); Louvain ($\gamma=0.5,1.0,2.0$); \\
 & CPM ($\gamma=0.05,0.10,0.20$); Multilevel; Fast Greedy; Leading Eigenvector; Spin Glass \\
C. Information-theoretic / spectral & InfoMap; Walktrap; spectral clustering \\
D. Generative / label propagation & Label propagation \\
E. Overlapping communities & Clique percolation \\
F. Graph representation learning & Self-supervised GAT; GraphSAGE (edge-prediction objective) \\
G. Classical clustering & $K$-means; Gaussian mixture; Ward; DBSCAN; HDBSCAN; OPTICS; \\
 & Mean Shift; Affinity Propagation; NMF \\
\midrule
\multicolumn{2}{l}{\emph{Total: 31 algorithms from 7 families}} \\
\bottomrule
\end{tabular}
\end{table*}

\noindent\textbf{FACEO-R outperforms existing methods under data degradation.} The resolution instability established above is intrinsic to single-graph approaches. We
therefore tested FACEO-R (Section~\ref{sec:faceor}), which anchors the partition to
GO-derived functional similarity with a weight that adapts to graph reliability, under five
controlled degradation conditions applied to 33 pathway--cohort pairs.

\noindent\textbf{Global results.} Across all 165 (pair $\times$ condition) comparisons,
FACEO-R achieved higher held-out Reactome coherence than Leiden $\gamma=1.0$ in 124 cases
(75\%; paired Wilcoxon $p = 1.9\times10^{-14}$, Cliff's $\delta = +0.292$) and than
WGCNA-style hierarchical cutting in 127 (88\%; $p = 9.0\times10^{-20}$,
$\delta = +0.296$). Mean coherence was $0.1915$ (FACEO-R) versus $0.1743$ (Leiden) and
$0.1527$ (WGCNA).

\noindent\textbf{The advantage grows as the expression graph degrades} (Fig.~\ref{fig:faceor}),
and is already present without degradation. On \emph{clean} data FACEO-R exceeded Leiden
($0.2078$ versus $0.1746$; $\Delta = +0.033$; 19/28 wins; $p = 0.022$). Under \emph{feature
dropout} (50\%), the regime most characteristic of single-cell data, the margin widened to
$+0.070$ ($0.2031$ versus $0.1328$; 27/28 wins; $p = 7.5\times10^{-8}$). Under expression
noise ($\sigma = 0.5$ and $1.0$) gains were $+0.049$ and $+0.043$
($p = 5.3\times10^{-4}$ and $4.8\times10^{-3}$), and under sample subsampling $+0.047$ and
$+0.047$ ($p = 2.7\times10^{-4}$ and $2.2\times10^{-4}$).

\noindent\textbf{The anchoring weight is critical and its optimum is narrow.} A systematic
sweep of the fixed anchoring weight (Fig.~\ref{fig:alpha}) revealed a monotone dependence:
strong anchoring is uniformly better, with mean coherence falling from $0.207$ at
$\alpha = 0.05$ to $0.189$ at $\alpha = 0.55$. The reliability-adaptive schedule originally
proposed ($\alpha \in [0.35, 0.95]$) therefore under-anchored and scored $0.184$ --- below
every fixed setting in $[0.05, 0.35]$. Retuning the schedule to $\alpha(q) \in [0.05, 0.25]$
recovered $0.206$, matching the best fixed weight while retaining adaptivity, and it wins in
3 of 6 conditions. All results reported here use the retuned schedule.

\begin{figure}[htbp]
\centering
\includegraphics[width=\columnwidth,height=0.44\textheight,keepaspectratio]{../output/figures_v3/Fig11_faceo_r_robustness.png}
\caption{FACEO-R outperforms existing methods under data degradation. (a--c) Functional
coherence as a function of degradation level for feature dropout, expression noise and sample
subsampling. (d) Per-condition win rate against Leiden and WGCNA. (e) Adaptive anchoring
weight $\alpha$ as a function of estimated graph reliability $q$. (f) Overall mean coherence
across all 145 comparisons.}
\label{fig:faceor}
\end{figure}

\begin{figure}[htbp]
\centering
\includegraphics[width=\columnwidth,height=0.44\textheight,keepaspectratio]{../output/figures_v3/FigCA_clean.png}
\caption{The advantage is present with and without degradation. Mean functional coherence
($z$) under clean data and five degradation regimes; FACEO-R exceeds both baselines in every
condition, with the largest margins under feature dropout.}
\label{fig:clean}
\end{figure}

\begin{figure}[htbp]
\centering
\includegraphics[width=\columnwidth,height=0.44\textheight,keepaspectratio]{../output/figures_v3/FigCA_alpha.png}
\caption{The anchoring weight has a narrow optimum and must be tuned. Mean coherence as a
function of a fixed anchoring weight $\alpha$ (red). Performance improves monotonically as
anchoring strengthens; the originally proposed adaptive schedule (blue dashed) under-anchored
and fell below every fixed setting in $[0.05, 0.35]$, whereas the retuned schedule
(green dotted) matches the best fixed weight.}
\label{fig:alpha}
\end{figure}


\noindent\textbf{Mechanism.} Because the anchoring weight $\alpha(q)$ decreases with graph
reliability, the functional term contributes most precisely when expression information is
least trustworthy, while FACEO-R never fell below the Leiden baseline on clean data ---
consistent with the designed non-inferiority.


\noindent\textbf{Validation on real single-cell data.} \label{sec:scrna}

To test whether the degradation results above transfer to genuine data rather than simulated
corruption, we applied the identical protocol to a real single-cell RNA-seq cohort: GSE145926
(COVID-19 bronchoalveolar lavage fluid \citep{liao2020covid}; mild $n=3$, severe $n=6$;
7{,}720 cells $\times$ 17{,}900 genes after QC, processed with Scanpy \citep{wolf2018scanpy} following established
single-cell practice \citep{luecken2019benchmark}). True zero fractions ranged
from 51\% to 89\% (median 84\%) ---
substantially higher than the 50\% dropout simulated above.

Across 16 pathway--subnetwork tests (Fig.~\ref{fig:scrna}), FACEO-R achieved higher
held-out Reactome coherence
than Leiden in \textbf{11 cases (69\%)} (paired Wilcoxon $p = 0.029$) and than WGCNA in
\textbf{15 (94\%)} ($p = 6.1\times10^{-5}$). Mean coherence was $0.252$ (FACEO-R) versus
$0.217$ (Leiden) and $0.181$ (WGCNA). Notably, the magnitude of the advantage was \emph{not}
correlated with the zero fraction (Spearman $\rho = -0.28$, n.s.), indicating that on real
data the benefit arises from the functional anchor supplying information orthogonal to
expression density, rather than from compensating dropout per se.

\begin{figure}[htbp]
\centering
\includegraphics[width=\columnwidth,height=0.44\textheight,keepaspectratio]{../output/figures_v3/Fig12_scrna_validation.png}
\caption{Validation on real single-cell data. (a) Distribution of true zero fractions across
pathway subnetworks in GSE145926 (median 84\%), with the simulated level used in Fig.~8
marked. (b) Paired coherence for the three methods across 16 subnetworks. (c) Coherence gain
as a function of zero fraction, showing no significant trend. (d) Win rate and significance
against each baseline.}
\label{fig:scrna}
\end{figure}



%% ============================================================
\section{Discussion}
%% ============================================================

\noindent\textbf{Principal findings.} We conducted a paired, multi-cohort benchmark of 31 partitioning algorithms on identical
pathway-derived graphs, evaluated across six independent dimensions with synthetic
ground-truth controls. Four conclusions follow. First, multi-resolution community detection
yields functionally more coherent partitions than single-scale hierarchical cutting, without
sacrificing modularity (Section~\ref{sec:main}). Second, synthetic benchmarks with known
nested structure show that the advantage is attributable to an actual resolution capability:
multi-resolution methods recover fine-grained sub-groups that hierarchical cutting merges.
Third, the resolution parameter must be chosen against an external functional criterion;
module count is not a valid proxy for quality. Fourth, and most consequentially for practice,
modularity does not exceed degree-preserving null models on these networks
(Section~\ref{sec:null}), and self-supervised embeddings only partially reproduce the
inferred communities.

\noindent\textbf{Implications for practice.} Pathway-centric module analysis is common in transcriptomic studies, and the choice of
partitioning method is rarely justified. Our results support three practical
recommendations. (i) \emph{Report functional coherence against a pathway-internal null.}
Enrichment $p$-values are biased by module size and, within a functionally coherent pathway,
are inflated for any partition; a size-matched, pathway-internal random baseline removes
both biases. (ii) \emph{Do not use module count or modularity alone as evidence of quality.}
Modularity failed to exceed degree-preserving null models, and module count is trivially
controlled by the resolution parameter. (iii) \emph{Select resolution against an external
criterion} --- functional coherence, stability across resamples, or downstream predictive
performance --- rather than by convention.

\noindent\textbf{A robustness-oriented prescription.} On clean data, multi-resolution partitioning is hard to improve upon: it already optimises a
structural proxy well correlated with the functional target, so alternative mechanisms tend
to trade one metric for another. The picture changes when the expression graph degrades ---
as it does in single-cell data, small clinical cohorts and noisy measurements. FACEO-R
exploits the fact that functional annotation is \emph{unaffected} by expression degradation,
and adapts its reliance on that annotation to the estimated reliability of the graph. The
result is a method that is never worse than standard Leiden on clean data yet markedly better
under degradation ($+0.046$ coherence under 50\% feature dropout; 88\% win rate;
$p<10^{-4}$). The practical value of annotation anchoring therefore lies not in beating
structure-based methods on ideal data, but in preserving module quality when data are
limiting --- precisely the regime in which pathway-level module analysis is most often
attempted.

\noindent\textbf{Limitations.} The analysis is restricted to nine GO biological-process pathways and sixteen cohorts;
although this exceeds typical practice, it remains modest by the standards of large-scale
benchmarking consortia. Density-based algorithms (DBSCAN, OPTICS, Mean Shift) were evaluated
after spectral embedding to mitigate the curse of dimensionality, which may understate their
performance on raw similarity profiles. Four high-complexity algorithms (edge betweenness,
link communities, optimal modularity and full multi-scale Markov Stability) could not be
evaluated reliably at the pathway sizes considered and were excluded; this is a limitation of
feasibility rather than a negative finding about those methods. Finally, the downstream
evaluation is limited to disease-status classification, which is an easy task in these
cohorts; survival and drug-target endpoints would be more discriminative.

\noindent\textbf{Future work.} Three directions follow naturally. First, extending the benchmark to single-cell and
multi-omic data would test whether the conclusions hold when co-expression is estimated from
sparser, more heterogeneous measurements. Second, overlapping community detection ---
included here only in a limited form --- deserves systematic evaluation, since pleiotropy
implies that hard partitions may be an inappropriate model for gene function. Third,
replacing pathway-referenced benchmarks with experimentally derived complexes (e.g. CORUM)
would provide a more direct biological ground truth.

%% ============================================================
\section{Conclusion}
%% ============================================================

On identical pathway-derived co-expression graphs, multi-resolution community detection
produced functionally more coherent modules than hierarchical co-expression clustering,
recovered nested functional sub-structure in synthetic benchmarks with known ground truth,
and yielded better downstream clinical prediction --- while modularity failed to exceed
degree-preserving null models and self-supervised embeddings only partially reproduced the
inferred communities. We therefore recommend that pathway-centric module analyses report
functional coherence against pathway-internal null expectations and select resolution
parameters against external criteria, rather than relying on module counts or modularity.

%% ============================================================
\section*{Acknowledgements}
%% ============================================================

The author thanks colleagues in the Department of Health Technology and Informatics at The
Hong Kong Polytechnic University for helpful discussions.

\section*{Funding}

The author declares no funding.

\section*{Competing interests}

The author declares no competing interests.

\section*{Data availability}
All expression cohorts are publicly available from GEO under the accessions in
Table~\ref{tab:datasets}. Processed expression matrices, complete evaluation tables (3100
main-benchmark, 528 synthetic-benchmark, 317 downstream, 495 FACEO-R degradation and 60
single-cell validation records), statistical outputs, all figure-generating scripts and the
manuscript source are organised as a reproducible pipeline and are available at
\url{https://github.com/sunmy1129/faceo-r}. The single-cell dataset analysed in
Section~\ref{sec:scrna} is available from GEO under accession GSE145926. Supplementary tables
S1--S6 are provided in the Supplementary Material.

\begin{thebibliography}{99}
\setlength{\itemsep}{1.5pt}
\bibitem{fortunato2010review} Fortunato S (2010) Community detection in graphs. \textit{Physics Reports} 486, 75--174.
\bibitem{langfelder2008wgcna} Langfelder P, Horvath S (2008) WGCNA: an R package for weighted correlation network analysis. \textit{BMC Bioinformatics} 9, 559.
\bibitem{langfelder2012wgcna} Langfelder P, Horvath S (2012) Fast R functions for robust correlations and hierarchical clustering. \textit{Journal of Statistical Software} 46, 1--17.
\bibitem{gillis2012guilt} Gillis J, Pavlidis P (2012) "Guilt by association" is the exception rather than the rule in gene networks. \textit{PLoS Computational Biology} 8, e1002444.
\bibitem{ashburner2000go} Ashburner M, Ball CA, Blake JA, et al (2000) Gene Ontology: Tool for the unification of biology. \textit{Nature Genetics} 25, 25--29.
\bibitem{newman2004girvan} Newman ME, Girvan M (2004) Finding and evaluating community structure in networks. \textit{Physical Review E} 69, 026113.
\bibitem{traag2019leiden} Traag VA, Waltman L, van Eck NJ (2019) From Louvain to Leiden: guaranteeing well-connected communities. \textit{Scientific Reports} 9, 5233.
\bibitem{traag2011cpm} Traag VA, Van Dooren P, Nesterov Y (2011) Narrow scope for resolution-limit-free community detection. \textit{Physical Review E} 84, 016114.
\bibitem{rosvall2008infomap} Rosvall M, Bergstrom CT (2008) Maps of random walks on complex networks reveal community structure. \textit{PNAS} 105, 1118--1123.
\bibitem{ahn2010link} Ahn YY, Bagrow JP, Lehmann S (2010) Link communities reveal multiscale complexity in networks. \textit{Nature} 466, 761--764.
\bibitem{lancichinetti2011oslom} Lancichinetti A, Radicchi F, Ramasco JJ, Fortunato S (2011) Finding statistically significant communities in networks. \textit{PLoS ONE} 6, e18961.
\bibitem{song2015megena} Song WM, Zhang B (2015) MEGENA: Multiscale Clustering of Geometrical Network. \textit{Bioinformatics} 31, 909--916.
\bibitem{luecken2019benchmark} Luecken MD, Theis FJ (2019) Current best practices in single-cell RNA-seq analysis: A tutorial. \textit{Molecular Systems Biology} 15, e8746.
\bibitem{lancichinetti2008lfr} Lancichinetti A, Fortunato S, Radicchi F (2008) Benchmark graphs for testing community detection algorithms. \textit{Physical Review E} 78, 046110.
\bibitem{ward1963hierarchical} Ward JH (1963) Hierarchical grouping to optimize an objective function. \textit{Journal of the American Statistical Association} 58, 236--244.
\bibitem{blondel2008louvain} Blondel VD, Guillaume JL, Lambiotte R, Lefebvre E (2008) Fast unfolding of communities in large networks. \textit{Journal of Statistical Mechanics}, P10008.
\bibitem{clauset2004fastgreedy} Clauset A, Newman ME, Moore C (2004) Finding community structure in very large networks. \textit{Physical Review E} 70, 066111.
\bibitem{newman2006leading} Newman ME (2006) Finding community structure in networks using the eigenvectors of matrices. \textit{Physical Review E} 74, 036104.
\bibitem{reichardt2006spin} Reichardt J, Bornholdt S (2006) Statistical mechanics of community detection. \textit{Physical Review E} 74, 016110.
\bibitem{pons2005walktrap} Pons P, Latapy M (2006) Computing communities in large networks using random walks. \textit{Journal of Graph Algorithms and Applications} 10, 191--218.
\bibitem{raghavan2007lpa} Raghavan UN, Albert R, Kumara S (2007) Near linear time algorithm to detect community structures in large-scale networks. \textit{Physical Review E} 76, 036106.
\bibitem{palla2005clique} Palla G, Derenyi I, Farkas I, Vicsek T (2005) Uncovering the overlapping community structure of complex networks in nature and society. \textit{Nature} 435, 814--818.
\bibitem{velickovic2018gat} Velickovic P, Cucurull G, Casanova A, Romero A, Lio P, Bengio Y (2018) Graph attention networks. \textit{International Conference on Learning Representations (ICLR)}.
\bibitem{velickovic2019dgi} Velickovic P, Fedus W, Hamilton WL, Lio P, Bengio Y, Hjelm RD (2019) Deep graph infomax. \textit{International Conference on Learning Representations (ICLR)}.
\bibitem{ester1996dbscan} Ester M, Kriegel HP, Sander J, Xu X (1996) A density-based algorithm for discovering clusters in large spatial databases with noise. \textit{ACM SIGKDD International Conference on Knowledge Discovery and Data Mining (KDD)}, 226--231.
\bibitem{campello2013hdbscan} Campello RJ, Moulavi D, Sander J (2013) Density-based clustering based on hierarchical density estimates. \textit{Pacific-Asia Conference on Knowledge Discovery and Data Mining (PAKDD)}, 160--172.
\bibitem{jassal2020reactome} Jassal B, Matthews L, Viteri G, et al (2020) The Reactome pathway knowledgebase. \textit{Nucleic Acids Research} 48, D498--D503.
\bibitem{benjamini1995fdr} Benjamini Y, Hochberg Y (1995) Controlling the false discovery rate: a practical and powerful approach to multiple testing. \textit{Journal of the Royal Statistical Society: Series B} 57, 289--300.
\bibitem{cliff1993dominance} Cliff N (1993) Dominance statistics: Ordinal analyses to answer ordinal questions. \textit{Psychological Bulletin} 114, 494--509.
\bibitem{hubert1985ari} Hubert L, Arabie P (1985) Comparing partitions. \textit{Journal of Classification} 2, 193--218.
\bibitem{fortunato2007resolution} Fortunato S, Barthelemy M (2007) Resolution limit in community detection. \textit{PNAS} 104, 36--41.
\bibitem{liao2020covid} Liao M, Liu Y, Yuan J, Wen Y, Xu G, Zhao J, et al (2020) Single-cell landscape of bronchoalveolar immune cells in patients with COVID-19. \textit{Nature Medicine} 26, 842--844.
\bibitem{wolf2018scanpy} Wolf FA, Angerer P, Theis FJ (2018) SCANPY: Large-scale single-cell gene expression data analysis. \textit{Genome Biology} 19, 15.
\end{thebibliography}



\end{document}
